\documentclass{article}

\usepackage{../../Style}

\title{Symplectic topology}
\author{Shing Tak Lam}
\date{Lent 2023\footnote{Based on lectures by Abigail Ward. Last updated \today.}}

\makeindex

\begin{document}

\maketitle

\tableofcontents

\lec

\section{Motivation}

Symplectic topology is the study of global phenomenon of \emph{symplectic manifolds}, which are a pair \((M, \omega)\), where \(M\) is a smooth manifold, and \(\omega\) is a non-degenerate closed \(2\)-form on \(M\). In particular, non-generacy implies the map
\begin{align*}
    \T M &\to \T^* M \\
    v &\mapsto \omega(v, \cdot)
\end{align*}
is an isomorphism.

\begin{example}
    \(\R^{2n}\) with coordinates \(x_1, y_1, \dots, x_n, y_n\) is a symplectic manifold, with
    \[\omega_{\text{std}} = \omega_0 = \sum \dd x_i \wedge \dd y_i\]
\end{example}

\begin{example}
    Let \(\Sigma\) be a surface of genus \(g\). Then any volume form makes \(\Sigma_g\) into a symplectic manifold.
\end{example}

Why study such manifolds? Some motivation from classical mechanics: Say we have a phase space \(X\) encoding position and momentum of some particles in a physical system. For example, we have a pendulum with angular coordinate \(\theta\), and angular momentum \(p = m\dot\theta\). The phase space is then \(S^1 \times \R\).

Given \(H : X \to \R\), encoding the total energy of a configuration, we want a vector field \(\mcX_H\) such that the trajectory \(x(t)\) satisfies
\[\dot x(t) = \mcX_H(x(t))\]
It's not enough to know that \(H\) is constant on trajectories, i.e. \(\dd H(\mcX_H) = 0\), since this is a co-dimension 1 condition. Recall if \(X\) has symplectic form \(\omega\), then we can take \(\mcX_H\) such that
\begin{equation}
    \label{eq:ham-vf}
    \omega(\mcX_H, \cdot) = \dd H
\end{equation}
Indeed, every phase space has such a symplectic form. Note that \(\omega\) is non-degenerate if and only if \cref{eq:ham-vf} has a solution for all \(H\), and \(\omega\) being closed is that \(\omega\) is preserved under the flow of \(\mcX_H\).

\begin{example}
    [harmonic oscillator] In this case, \(X = \R^2\), with coordinates \(p\) for momentum, \(q\) for position. The symplectic form is
    \[\dd p \wedge \dd q = r \dd r \wedge \dd\theta\]
    where \((r, \theta)\) is polar coordinates on \(X\).
    This has Hamiltonian
    \[H = \frac{1}{2}(p^2 + q^2) = \frac{1}{2}r^2\]
    Hence
    \[\dd H = r \dd r\]
    Thus, \(\mcX_H = \pdv{\theta}\).
\end{example}

\subsection{Mathematical perspective}

A question to keep in mind is: How does symplectic geometry compare to e.g. Riemannian geometry, complex geometry etc.

In this context, \(\dd\omega = 0\) is a form of ``flatness''. Suppose \((Y, g)\) is a Riemannian manifold, i.e. we have \(g \in \Sym^2\T^* Y\) non-degenerate. It is a fact that \((Y, g)\) is locally isometric to \(\R^n\) with the Euclidean metric if and only if the curvature \(F\) of \(g\) vanishes. That is, \(g\) is flat. 

On the other hand, suppose \((Z, J)\) is an almost complex manifold. That is, \(J \in \End(\T Z)\), with \(J^2 = - \id_{\T Z}\). One example would be \(\C^n\), with \(J(z) = iz\). One question is, when is \((Z, J)\) locally isomorphic to \(\C^n\)? That is, when is \((Z, J)\) a complex manifold. A difficult result in analysis, called the Newlander-Nirenberg theorem, shows that this occurs if and only if the Nijenhuis tensor,
\[\mcN_J(v, w) = [Jv, Jw] - J([v, Jw] + [Jv, w]) - [v, w]\]
vanishes

\begin{theorem}
    [Darboux] A manifold \((X, \omega)\) with non-degenerate \(2\)-form \(\omega\), is locally modeled on \((\R^{2n}, \omega_{std})\) if and only if \(\dd\omega = 0\).
\end{theorem}

Finally, symplectic topology is ubiquitous in research in geometry and topology, e.g. global topology, algebraic geometry and so on.

\lec

\section{(Symplectic) linear algebra}

Let \(V\) be a vector space, over a field \(\mathbb F\) of charactistic 0, with \(\dim(V) = n\). \(V^* = \Hom(V, \mathbb F)\) will denote the dual.

Define \(\mathbb F\)-algebras
\begin{align*}
    \T(V) &= \bigoplus_{i=0}^\infty V^{\otimes i} \\
    \Lambda^*(V) &= \frac{\T(V)}{\inner{v \otimes v}}
\end{align*}
If \(\{e_1, \dots, e_n\}\) is a basis for \(V\), then
\[\{e_{i_1} \wedge \cdots \wedge e_{i_k} \mid i_1 < \cdots < i_k\}\]
is a basis for \(\Lambda^kV\). Using this,
\[\dim(\Lambda^kV) = {n \choose k}\]

\begin{proposition}
    \(\Lambda^k(V^*)\) is isomorphic to the space of alternating multilinear maps
    \[V^{\otimes k} \to \bb F\]
\end{proposition}

\begin{proof}
    The map is given by
    \[e_{i_1}^\vee\wedge \cdots \wedge e_{i_k}^\vee \mapsto \sum_{\sigma \in S_k}\sign(\sigma)e_{i_{\sigma(1)}}^\vee\otimes \cdots \otimes e_{i_{\sigma(k)}}^\vee\]
    The rest is left as an exercise.
\end{proof}

In particular, \(\Lambda^n(V^*) = \Lambda^{\text{top}}(V^*)\) is (non-canonically) isomorphic to \(\Span\{\det\}\), or volume forms.

\begin{definition}
    A \emph{symplectic form} on \(V\) is \(\omega \in \Lambda^2(V^*)\), such that \(\omega : V \times V \to \bb F\) is a perfect pairing\footnote{That is, \(\omega(v, \cdot)\) induces an isomorphism \(V \cong V^*\).}.
\end{definition}

\begin{example}
    Let \(M^{4m+2}\) be an oriented manifold, and \(V = \HH^{2m+1}(M; \Q)\). The pairing
    \[\inner{\alpha, \beta} = \int_M\alpha\cup\beta\]
    is a symplectic form.
\end{example}

\begin{example}
    Let \(V = \R^{2n}\), with basis \(\{x_i, y_i\}\), with dual basis \(\{e_i, f_i\}\). We can define
    \[\omega = \sum_{i=1}^n e_i \wedge f_i\]
    More generally, if \((V, \omega_V), (W, \omega_W)\) are symplectic, then so is
    \[(V \oplus W, \omega_V \oplus \omega_W)\]
\end{example}

We can write
\[\omega(v, w) = v \vdot J_0 w\]
where
\[J_0e_j = f_j \qqtext{and} J_0f_j = -e_j\]
If we use complex coordinates \(z_j = x_j + iy_j\), then \(J_0\) is just multiplication bu \(i\).

Let \((V, \omega)\) be a symplectic vector space. Let \(W \sse V\) be a subspace. Define
\[W^\omega = \{v \in V \mid \omega(w, v) = 0 \text{ for all }w \in W\}\]

\begin{definition}
    \(W\) is
    \begin{enumerate}
        \item \emph{isotropic} if \(W \sse W^\omega\),
        \item \emph{coisotropic} if \(W^\omega \sse W\),
        \item \emph{symplectic} if \(W \cap W^\omega = 0\),
        \item \emph{Lagrangian} if \(W = W^\omega\). That is, it is isotropic and coisotropic.
    \end{enumerate}
\end{definition}

\begin{example}
    Take \(\R^{2n}\) with the standard symplectic form
    \[\omega = \sum_i \dd x_i \wedge \dd y_i\]
    Then \(\Span\{x_1, x_2\}\) is isotropic. In fact, \(\Span\{x_1, \dots, x_n\}\) is a Lagrangian subspace. On the other hand, \(\Span\{x_1, y_1\}\) is a symplectic subspace.
\end{example}

\begin{remark}
    For \emph{any} subspace \(W \sse V\),
    \[\dim(W) + \dim(W^\omega) = \dim(V)\]
    and we have that \((W^\omega)^\omega = W\).
\end{remark}

\begin{theorem}
    Given any symplectic vector space \((V, \omega)\) over \(\R\), there exists a linear isomorphism \(\Phi : V \to \R^{2n}\), such that \(\Phi^*\omega_0 = \omega\).
\end{theorem}

\begin{proof}
    We induct on \(\dim(V)\). Choose any \(v_1 \in V\) non-zero. By non-degeneracy, there exists \(w_1 \in V\) such that \(\omega(v_1, w_1) = 1\). Let \(V_1 = \Span\{v_1, w_1\}\). Define \(\Phi_1 : V_1 \to \R^2\), by \(\Phi_1(v_1) = x_1\) and \(\Phi_1(w_1) = y_1\). This satisfies that \(\Phi_1^*\omega_0 = \omega\vert_{V_1}\).

    Define \(V_2 = V_1^\omega\). Then \(V_1\) and \(V_2\) are symplectic, and \(\dim(V_2) = \dim(V) - 2\). By induction hypothesis, there exists \(\Phi_2 : V_2 \to \R^{2n-2}\), with \(\Phi_2^*\omega_0 = \omega\vert_{V_2}\). We can then define
    \[\Phi = \Phi_1 \oplus \Phi_2 : (V_1, \omega\vert_{V_1}) \oplus (V_2, \omega\vert_{V_2}) \to (\R^{2n-2}, \omega_0) \oplus (\R^2, \omega_0) = (\R^{2n}, \omega_0)\]
\end{proof}

\lec

In particular, symplectic vector spaces are always even dimensional.

\begin{corollary}
    \(\alpha \in \Lambda^2 V^*\) is symplectic if and only if \(\alpha \wedge \cdots \wedge \alpha \in \Lambda^{\text{top}} V^* = \R\) is non-zero.
\end{corollary}

\begin{proof}
    Note that
    \[\omega_{\text{std}}^n = x_1^\vee \wedge y_1^\vee \cdots \wedge x_n^\vee \wedge y_n^\vee\]
    is the usual volume form.\footnote{or maybe with a coefficient of \(n!\) in front of it...}
\end{proof}

Returning to linear algebra, we have
\[\Hom(V, V) \to \Hom(\Lambda^\text{top}V, \Lambda^\text{top}V) \cong \R\]
The composition is the \emph{determinant} \(\det : \Hom_*(V, V) \to \R\).

\subsection{Orientations}

Define
\[\ori_V = \frac{\{\text{ordered bases for }V\}}{\GL_+(V)} = \frac{\Fr(V)}{\GL_+(V)}\]
where we consider the \(\GL_+(V)\)-orbits, i.e. two frames are equivalent if they are related by a transformation with positive determinant.

In this case, \(\ori_V\) has two elements. Given any volume form \(\alpha\) on \(V\), there exists an induced orientation on \(V\), given by the sign of \(\alpha\) on a basis. In particular, any symplectic vector space has such a canonical orientation, induced by \(\omega^n\).

\subsection{Symplectomorphism}

\begin{definition}
    Define
    \[\Sp(2n) = \{\Phi \in \GL(\R^{2n}) \mid \Phi^*\omega_{\std} = \omega_{\std}\}\]
    if \(J_0\) is the matrix so that \(\omega(u, v) = u \vdot J_0 v\), then we have that
    \[\Sp(2n) = \{\Phi \in \GL(\R^{2n}) \mid \Phi^\T J_0\Phi = J_0\}\]
\end{definition}

\begin{example}
    For \(\R^2\), \(\omega\) is the standard volume form \(\dd x \wedge \dd y\). So preserving \(\omega\) is the same as preserving the volume form. Hence
    \[\Sp(2) = \SL(2, \R)\]
    In general, we have an embedding
    \[\Sp(2n) \sse \SL(2n, \R)\]
    since \(\omega^n\) is the standard volume form\footnote{maybe up to a factor of \(n!\)} on \(\R^{2n}\).
\end{example}

Note if \(\Phi\) commutes with \(J_0\), and \(\Phi\) preserves the inner product, then \(\Phi \in \Sp(2n)\). So
\[\GL(n, \C) \cap \mathrm O(2n) \sse \Sp(2n)\]
In fact, we have
\[\GL(n, \C) \cap \mathrm O(2n) = \GL(n, \C) \cap \Sp(2n) = \Sp(2n) \cap \mathrm O(2n) = \mathrm U(n)\]

\begin{proposition}
    The set \(S\) of symplectic forms on \(\R^{2n}\), as a subset of \(\Lambda^2(V^*)\) is \(\GL(2n, \R)/\Sp(2n)\), a homoegenous space.
\end{proposition}

\begin{proof}
    \(\GL(2n, \R)\) acts transitively on \(S\), with stabiliser \(\Sp(2n)\) for any point.
\end{proof}

Let \(\Lag(V)\) be the space of Lagrangian subspaces \(L\) of \(V\).

\begin{example}
    \(\Lag(\R^2)\) is just the one-dimensional subspaces of \(\R^2\), i.e. \(\R\P^1 = S^1\). More generally,
    \[\Lag(\R^{2n}) = \frac{\mathrm U(n)}{\mathrm O(n)}\]
\end{example}

\begin{example}
    Let \(V = C^\infty(S^1, \R)\), with
    \[\omega(f, g) = \int_{S^1}f\dd g\]
    Using integration by parts, this is alternating.
\end{example}

\subsection{Vector bundles}

If \(\pi : \mcE \to X\) is a vector bundle, we can define \(\Lambda^*\mcE \to X\), \(\Lambda^*(\mcE^*) \to X\), \(\ori_\mcE \to X\), and so on. We have a vector space \(\Gamma(\mcE)\) of global sections of \(\mcE\). In particular,
\[\Omega^k(X) = \Gamma(\Lambda^k(\T^* X))\]
is the space of \(k\)-forms.

\lec

\begin{definition}
    \(\mcE\) is \emph{orientable} if there exists a global section of \(\ori_\mcE\).
\end{definition}

\begin{example}
    [non-example] The M\"obius bundle over \(S^1\) is not orientable. More generally, for any vector bundle, \(\mcE \to S^1\), \(\mcE\) is trivial if and only if it is orientable.
\end{example}

Let \(M\) be a (connected smooth) manifold, and consider the tangent bundle \(\T M \to M\).

\begin{definition}
    \(M\) is \emph{orientable} if \(\T M\) is. If so, a choice of orientation for \(M\) is a choice of orientation for \(\T M\).
\end{definition}

Note that if \(M\) is orientable, then there exists two choices of orientation. If \(M\) represents one orientation, we will write \(\overline{M}\) for the other. Moreover, if \(M\) has a boundary, then \(\partial M\) has a canonical orientation, by saying \((e_1, \dots, e_{n-1})\) is positively oriented if \((\nu, e_1, \dots, e_{n-1})\) is positively oriented, where \(\nu\) is the outward pointing normal.

For compact \(M\), \(M\) is orientable if and only if \(\HH^n(M; \Z) \cong \Z\). A choice of orientation corresponds to a choice of generator in \(\{\pm 1\}\). Note that elements of \(\Diff(M)\) must either preserve or reverse orientation, which we can read off from \(\varphi_* : \HH^n(M; \Z) \to \HH^n(M; \Z)\).

\subsection{Cohomology and integration}

\begin{theorem}
    [de Rham] We have an isomorphism between de Rham cohomology and singular homology. Moreover, this is an isomorphism of rings, between the wedge product and the cup product.
\end{theorem}

Intuitively, the isomorphism is given by integrating a form over a chain.

\begin{theorem}
    [Stokes]
    Let \(M\) be a smooth \(n\)-manifold with boundary, and \(\alpha\) a compactly supported \((n-1)\)-form on \(M\). Then
    \[\int_M \dd\alpha = \int_{\partial M}\alpha\]
\end{theorem}

Recall that if \((M^{2n}, \omega)\) is symplectic, then \(\HH^{2k}(M; \R) \ne 0\) for all \(0 \le k \le n\).

\printindex

\end{document}